\documentclass[11pt]{article}
\usepackage{amsmath,amssymb,amsthm,hyperref}
\newtheorem{proposition}{Proposition}
\newcommand{\E}{\mathbb E}
\title{An explicit original-law entropy dual and a stronger necessary threshold}
\author{Campaign 3395: Engineer continuation}
\date{November 5, 2026}
\begin{document}
\maketitle

\paragraph{Status against the requested upper bound.}
CDR18 and MP9 remain unproved. This note gives an explicit challenger
to that upper bound under the entire unchanged original $\Delta$,
not another finite calibration. It also improves the internal RI7
information lower bound from $33N/2^{143}-4$ to
$37N/2^{143}-5/4$ at the same native loss and admissible threshold.
This is a stronger necessary condition on a retained cover, not
the requested original-DAG construction or circuit lower bound.
No achieved $K_0,a,\beta,b_*,N_0$, positive defect threshold, fixed
rational factor, verified novelty or publishable significance is
asserted. DS16--DS18 remain the strongest achieved unrestricted
circuit result; the six-fifths and four-thirds claims are unchanged.

\paragraph{What is reused rather than rerun.}
RI1 supplies the original parity law and RI4 the exact selector
identity. BP1 gives $R/\kappa\le h_Q\le\kappa R$ on all one-hot
assignments for $1<\kappa\le65/64$. The exact endpoint $\kappa=1$
can also be included. The original discrepancy theorem gives
\[
 |\E_\Delta\chi1_B|\le\gamma,\qquad
 \gamma=3\exp(-N/2^{136})
\]
for every balanced original-source rectangle, with
$\chi=1$ on non-trees and $-1$ on trees.
These reviewed statements are read in
\path{rectangle-selector-information.tex},
\path{all-box-pinning.tex}, and
\path{../sparse-sensitive-spanning-tree/theorem.tex}.
ED2 and CDR18 retain their meanings.
The previously documented repeated-literal-power issue in L10--L11
was found before a new computation; no duplicate power diagnostic
or quotient operation was executed.
No new external theorem is a premise.

\section{A scalar majorant that uses all rectangle masses}
Let $0<\gamma<1$ and put $\ell=\log(1/\gamma)>0$.
Convexity of $\exp$ on $[0,\ell]$ gives
\[
 \min\{1,\gamma e^r\}
       \le \gamma+(1-\gamma)\frac r\ell
                         \quad(r\ge0).               \tag{SD1}
\]
For $0\le r\le\ell$, the right side is the secant joining the
two endpoint values of $\gamma e^r$.
For $r\ge\ell$ it is at least one.
This is a pointwise inequality on the entire half-line, not a
numerically sampled comparison.
It replaces the RI12 split into rectangles above and below
$128\gamma$; no sampler, source partition or cover is changed.

For any nonnegative rectangle sum $P=\sum_Bw_B1_B$ of positive
mean $Z$, let $\pi_B=w_Bp(B)/Z$ with $p=\Delta$.
Only terms of positive $\pi_B$ occur in these expectations; all
full-domain domination constraints, including zero-law points,
remain in the definition of the cover. Discrepancy and (SD1) imply
\[
 a_P:=\frac{\E_p(\chi P)}Z
 \le\E_{\pi_B}\min\{1,\gamma/p(B)\}
 \le\gamma+\frac{1-\gamma}{\ell}
                     \frac{C_p(w)}Z.
\]
Consequently
\[
             \frac{C_p(w)}Z
               \ge\frac{\ell}{1-\gamma}(a_P-\gamma).
                                                               \tag{SD2}
\]
The two sides of a rectangle are not made independent under $p$.
The entire original probability of each rectangle is used.

\section{Optimize the native correlation and density bounds}
Assume the original full-domain domination $0\le P\le h_Q$ and
the single total loss $\E_p(h_Q-P)\le\epsilon=1/16$.
Write
$A_h=\E_p(h_Q1_{\{\chi=1\}})$ and
$B_h=\E_p(h_Q1_{\{\chi=-1\}})$.
The reviewed parity bias $\delta=2^{-q}\ge0$ gives
\[
 A_h\ge(1+\delta)/\kappa,\qquad
 B_h\le\kappa(1-\delta)/2.
\]
Thus $A_P\ge1/\kappa-\epsilon>0$ and $B_P\le\kappa/2$.
The ratio $(A-B)/(A+B)$ increases with $A>0$ and decreases
with $B\ge0$. Since $1/\kappa-\epsilon$ decreases and $\kappa/2$
increases with $\kappa$,
\[
 a_P\ge
 \frac{1/\kappa-\kappa/2-\epsilon}
      {1/\kappa+\kappa/2-\epsilon}
 \ge \rho:=\frac{3447}{11897}>\frac{37}{128}.         \tag{SD3}
\]
Unlike RI9, this retains the coupling between the numerator and
denominator rather than separately maximizing the total mass.
The same reasoning gives
\[
 \frac{\E_p(\chi h_Q)-\epsilon}{\E_ph_Q-\epsilon}
                              \ge\rho.               \tag{SD4}
\]
It is the same ratio with $A=A_h-\epsilon$ and $B=B_h$.
No new loss allowance is used in either calculation.

For clarity, multiplying the difference from $37/128$ in (SD3)
by the positive denominator gives, up to a positive constant,
\[
 182-165\kappa^2-\frac{91}8\kappa
 =\frac{45}8-\frac{2731}8t-165t^2
 \ge\frac{1027}{4096}>0,\quad
 t=\kappa-1\in[0,1/64].                              \tag{SD5}
\]
The polynomial is decreasing there and its endpoint value is
exact. The inspected interval operation independently encloses
the same whole interval.

The original target mean also improves the input-density bound:
\[
 Z\ge\frac3{2\kappa}-\frac1{16},\qquad
 \frac{P}{Z}\le
 \frac{4\kappa^2}{3-\kappa/8}
 \le\frac{4225}{2942}<\frac32 .
\]
Therefore, by the reviewed selector identity and $\log x\le x-1$,
\[
 0\le D(pP/Z\Vert p)
 \le\log(4225/2942)<\frac12.                          \tag{SD6}
\]
This bounds a density, not rectangle weight or circuit size.

\section{The improved all-size information diagnostic}
\begin{proposition}
Under the original BP1 promise, every full-domain dominated
balanced-original-source rectangle cover with total original-law
loss at most $1/16$ satisfies, for admissible $N\ge2^{142}$,
\[
 \frac{C_\Delta(w)}Z>\frac{37N}{2^{143}}-\frac34,
 \qquad
 I_\pi(\nu;B)>\frac{37N}{2^{143}}-\frac54.             \tag{SD7}
\]
The more precise first lower bound is
$\ell(\rho-\gamma)/(1-\gamma)$.
\end{proposition}
\begin{proof}
Put $x=N/2^{136}\ge64$, so $\ell=x-\log3$.
The exponential series gives $e^x\ge x^3/6$, and hence
\[
 \gamma\le18/x^3\le9/131072<1/2,\qquad
 \frac{\gamma\ell}{1-\gamma}
       \le2\gamma x\le36/x^2\le9/1024.
\]
Using $\rho<1/3$ and $\log3<2$,
\[
 \begin{split}
 \frac{\ell(\rho-\gamma)}{1-\gamma}
 &=\rho x-\rho\log3
                -(1-\rho)\frac{\gamma\ell}{1-\gamma}\\
 &>\rho x-\frac23-\frac9{1024}
 >\frac{37x}{128}-\frac34 .
 \end{split}
\]
The last comparison uses
$3/4-2/3-9/1024=229/3072>0$.
Apply (SD2)--(SD3) for the normalized cost.
Then subtract the input divergence in (SD6), using RI4, to obtain
the information bound. There is one cover and one loss budget.
\end{proof}

This improves RI7 by $N/2^{141}+11/4$ in its displayed right-hand
side. It does not improve DS16--DS18: their achieved circuit
upper-to-lower transfer is a different result, and the missing
MP9 upper estimate is not supplied by an information lower bound.
No finite calculation is extrapolated to establish (SD7).

\section{An explicit challenger to the native dual inequality}
\begin{proposition}
With the same original $p,\chi,\gamma,\ell$, let
$\lambda=\ell/(1-\gamma)$. For every $0\le T<\ell$, set
\[
 z_\nu=p_\nu\{\lambda(\chi(\nu)-\gamma)-T\},
 \qquad \eta=\ell-T>0.                               \tag{SD8}
\]
These are feasible in ED2 for the entire original balanced
rectangle family and all original domain coordinates. Moreover
\[
 z\cdot h_Q-\eta/16
 \ge(\E_ph_Q-1/16)
          \left\{\frac{\ell(\rho-\gamma)}{1-\gamma}-T\right\}.
                                                               \tag{SD9}
\]
In particular they violate the proposed native upper bound whenever
$0\le T\le37N/2^{143}-3/4$.
\end{proposition}
\begin{proof}
For each coordinate,
$z_\nu\le p_\nu(\lambda(1-\gamma)-T)=\eta p_\nu$.
At zero-law coordinates $z_\nu=0$, so their cap is also retained.
For a rectangle of probability $m>0$,
$\E_p\chi1_B\le\min\{m,\gamma\}$.
Apply (SD1) with $r=\log(1/m)$ to obtain
\[
 \begin{split}
 z(B)
 &\le m\{\lambda(\min(1,\gamma/m)-\gamma)-T\}\\
 &\le m\log(1/m)-Tm=d_T(B).
 \end{split}
\]
If $m=0$, both $z(B)$ and $d_T(B)$ are zero.
This verifies every dual constraint, not just sampled rectangles.

Write $Z_h=\E_ph_Q$ and $q_h=\E_p(\chi h_Q)$.
Since $\ell=\lambda(1-\gamma)$, the exact identity is
\[
 z\cdot h_Q-\eta\epsilon
 =(Z_h-\epsilon)
   \left\{\lambda\left(\frac{q_h-\epsilon}{Z_h-\epsilon}
                                   -\gamma\right)-T\right\}.
                                                               \tag{SD10}
\]
Use (SD4) and $Z_h-\epsilon>0$ to prove (SD9).
The strict scalar bound in the proof of (SD7) gives the final
positive objective when $T$ is at most the displayed rational
threshold. That threshold is below $\ell$, since $\rho<1$,
so (SD8) is defined throughout this asserted challenge range.
\end{proof}

There is no inconsistency with the existence of exact-target
circuits: their original sizes may be exponentially large.
For the native candidate
\[
 T=K_0+a\log(M_n+2)+\beta\log(s+1)+N/2^{139},
\]
the existence of MP9 would imply
\[
 K_0+a\log(M_n+2)+\beta\log(s+1)
                    >\frac{21N}{2^{143}}-\frac34.    \tag{SD11}
\]
Equation (SD11) is \emph{conditional on MP9}, not a circuit lower
bound now proved. It shows directly which challenge an actual
original-DAG upper argument must exclude by its size dependence.
The logarithmic overhead cannot be omitted when testing (SD8).

For example, at the bare linear price $T=x/8$ the challenger has
$\eta=7x/8-\log3$, between $27x/32$ and $7x/8$ for $x\ge64$.
Its objective is positive by (SD9), with price gap greater than
$21x/128-3/4$. This is an original-law multiplier of order $N$,
not a freely chosen favorable multiplier from a small calibration.
The omitted logarithmic overhead is exactly what an original-size
theorem must pay; no small BP1 circuit is constructed by this test.

\section{The upper argument still has no overlapping-product estimate}
The first candidate continuation was a gate-weighted context charge.
The repository already records why literal powers require divided
differences rather than ordinary-derivative occurrence counts,
so no unchanged failure was rerun.
The explicit full-law challenge (SD8) now tests the actual sign and
scale in CDR18 without imposing a stronger intermediate invariant.
It does not produce an upper estimate at overlapping products.

Substituting (SD8) into the desired bound requires an original-DAG
argument that controls the target's signed imbalance with the
logarithmic size allowance. A4/A7 still provide a joint nonnegative
multi-hot conflict identity, not this signed one-hot comparison.
Neither (SD1) nor finite primal attainment changes that fact.
No separate negative-completion estimate, subgate BP1 assumption,
formula unfolding, fixed-cut replacement, or restricted pooling
family is used.

Thus CDR18 for every $\eta$ remains open at positive defect.
There are no achieved native constants to instantiate GL10--GL13,
no LR3 materialization, and no new fixed rational approximation
factor. The diagnostic is not a refutation of MP9 with its original
size term. The broad dual route and the original-law raw-cover
alternative are not retired.

\paragraph{Exact replay and its scope.}
\path{secant_original_law_dual_check.py} verifies the new rational
correlation and density constants, the whole-interval polynomial
enclosure in (SD5), and the all-size scalar remainder inequalities
using exact arithmetic. The inspected Jacobian operation
\texttt{polynomial.box.enclosure.compute} returns the exact enclosure
$[1027/4096,45/8]$ for the displayed polynomial on $[0,1/64]$.
Its exact payload, typed output, server/protocol and adapter status
are preserved in dedicated JSON files.
The secant inequality, all-rectangle dual feasibility, and original
sampler use are justified by the proofs above, not finite sampler
enumeration. The certificate checks no circuit-size upper bound
and does not certify MP9.
\end{document}
